\documentclass[11pt,a4paper]{article}
\usepackage[french]{babel}
\usepackage[utf-8]{inputenc}
\usepackage[T1]{fontenc}
\usepackage[margin=2cm]{geometry}
\usepackage{graphicx}
\usepackage{listings}
\usepackage{xcolor}
\usepackage{hyperref}
\usepackage{booktabs}
\usepackage{array}
\usepackage{amsmath}
\usepackage{fancyhdr}

% Configuration listings
\lstset{
  language=SQL,
  basicstyle=\ttfamily\small,
  breaklines=true,
  frame=single,
  backgroundcolor=\color{gray!10},
  keywordstyle=\color{blue},
  commentstyle=\color{gray},
  stringstyle=\color{red},
}

\pagestyle{fancy}
\lhead{USPV Judo}
\chead{Dossier d'Analyse Technique}
\rhead{Mehdi LAFAY}
\rfoot{Page \thepage}

\title{\textbf{Dossier d'Analyse Technique}\\
{\large Site Web USPV Judo}}
\author{Mehdi LAFAY\\
BUT Informatique --- 2\textsuperscript{e} année\\
IUT Belfort-Montbéliard}
\date{Stage : 7 avril -- 5 juin 2026}

\begin{document}

\maketitle
\tableofcontents
\newpage

% =====================================================
\section{Introduction et Contexte}

\subsection{Objectif du projet}
Le site USPV Judo est la plateforme numérique de gestion et de promotion du club historique de judo de Pont-de-Roide-Vermondans. Le projet vise à :
\begin{itemize}
  \item Créer une présence web professionnelle et accessible
  \item Mettre en place un système d'inscription et de gestion des adhérents
  \item Permettre la gestion de l'agenda des séances et tournois
  \item Assurer la sécurité des données personnelles (RGPD)
\end{itemize}

\subsection{Périmètre technique}
\begin{tabular}{lp{10cm}}
\toprule
\textbf{Domaine} & \textbf{Technologie} \\
\midrule
Frontend & Nuxt 3 (Vue 3) + TailwindCSS \\
Backend & Node.js (Nitro runtime) + Express middleware \\
Base de données & PostgreSQL 14+ (24 tables) \\
Sécurité & JWT + bcrypt + middleware CSRF \\
Déploiement & VPS Linux + GitHub Actions CI/CD \\
Versioning & Git / GitHub \\
\bottomrule
\end{tabular}

% =====================================================
\section{Architecture Générale}

\subsection{Vue d'ensemble}
\begin{verbatim}
┌─────────────────────────────────────────┐
│         Navigateur (Client)             │
│         Nuxt 3 SPA (Vue 3)              │
│         - Pages publiques               │
│         - Formulaires réactifs          │
│         - Gestion d'état (Pinia)        │
└──────────────┬──────────────────────────┘
               │ HTTPS / REST API
┌──────────────▼──────────────────────────┐
│      Serveur Nitro (Node.js)            │
│      - Routes API (/api/...)            │
│      - Middleware (JWT, CORS, CSRF)     │
│      - Logique métier                   │
│      - Sessions utilisateur             │
└──────────────┬──────────────────────────┘
               │ TCP/5432
┌──────────────▼──────────────────────────┐
│       PostgreSQL                        │
│       - Familles (adhérents)            │
│       - Inscriptions                    │
│       - Sessions de judo                │
│       - Tournois                        │
└─────────────────────────────────────────┘
\end{verbatim}

\subsection{Modèle Client-Serveur}
\begin{itemize}
  \item \textbf{Client} : Application Nuxt 3 (SPA) exécutée dans le navigateur
  \item \textbf{API} : Endpoints RESTful exposés par le serveur Nitro
  \item \textbf{Base de données} : PostgreSQL centralisée, accessible uniquement par le serveur
\end{itemize}

% =====================================================
\section{Détail des Composants}

\subsection{Frontend (Nuxt 3 + Vue 3)}

\subsubsection{Structure}
\begin{lstlisting}[language=bash]
nuxt-app/
├── pages/                    # Pages publiques
│   ├── index.vue            # Accueil
│   ├── inscription.vue      # Formulaire d'inscription
│   └── planning.vue         # Calendrier séances/tournois
├── components/              # Composants réutilisables
│   ├── InscriptionForm.vue
│   ├── SessionCard.vue
│   └── TournamentTable.vue
├── composables/             # Logique métier réutilisable
│   ├── useAuth.ts           # Authentification
│   └── useFormValidation.ts # Validation formulaires
├── stores/                  # État global (Pinia)
│   ├── auth.ts              # Utilisateur courant
│   └── sessions.ts          # Séances/tournois
└── middleware/              # Middleware client
    └── auth.ts              # Garde routes protégées
\end{lstlisting}

\subsubsection{Réactivité et State Management}
\begin{itemize}
  \item \textbf{Pinia} : Store global pour authentification et données session
  \item \textbf{Composition API} : Logique réutilisable via composables
  \item \textbf{TailwindCSS} : Framework CSS utilitaire pour design responsive
\end{itemize}

\subsection{Backend (Nitro / Express)}

\subsubsection{Architecture}
\begin{lstlisting}[language=bash]
server/
├── api/                    # Endpoints API
│   ├── auth.ts            # Login / register
│   ├── inscriptions/       # CRUD inscriptions
│   ├── families/           # Gestion adhérents
│   ├── sessions/           # Planning séances
│   └── tournaments/        # Gestion tournois
├── middleware/            # Middleware global
│   ├── auth.ts            # Vérification JWT
│   ├── csrf.ts            # Protection CSRF
│   └── cors.ts            # CORS policy
├── utils/                 # Utilitaires
│   ├── db.ts              # Pool PostgreSQL
│   └── crypto.ts          # Hash bcrypt, JWT
└── types/                 # TypeScript types
    └── models.ts          # Interfaces DB
\end{lstlisting}

\subsubsection{Flux d'Authentification}
\begin{enumerate}
  \item Client envoie \texttt{POST /api/auth/login} avec identifiant + mot de passe
  \item Serveur requête \texttt{SELECT FROM famille WHERE email = ...}
  \item Vérification du hash bcrypt du mot de passe
  \item Génération JWT signé (payload : id, email, rôle)
  \item Client stocke le JWT en localStorage
  \item Requêtes suivantes incluent \texttt{Authorization: Bearer <JWT>}
  \item Middleware \texttt{auth.ts} valide la signature JWT à chaque requête
\end{enumerate}

% =====================================================
\section{Base de Données}

\subsection{Schéma relationnel}

Le schéma PostgreSQL compte 24 tables réparties en 4 domaines :

\begin{tabular}{lp{7cm}}
\toprule
\textbf{Domaine} & \textbf{Tables principales} \\
\midrule
Adhérents & \texttt{famille}, \texttt{representant}, \texttt{enfant}, \texttt{contact} \\
Inscriptions & \texttt{inscription}, \texttt{inscription\_detail}, \texttt{forfait} \\
Séances & \texttt{session}, \texttt{ceinture}, \texttt{salle} \\
Tournois & \texttt{tournoi}, \texttt{tournoi\_inscription}, \texttt{categorie} \\
\bottomrule
\end{tabular}

\subsection{Exemple : Table \texttt{famille}}

\begin{lstlisting}[language=SQL]
CREATE TABLE famille (
  id SERIAL PRIMARY KEY,
  email VARCHAR(255) UNIQUE NOT NULL,
  nom VARCHAR(100) NOT NULL,
  prenom VARCHAR(100),
  telephone VARCHAR(20),
  adresse TEXT,
  actif BOOLEAN DEFAULT true,
  password_hash VARCHAR(255) NOT NULL,
  created_at TIMESTAMP DEFAULT CURRENT_TIMESTAMP,
  updated_at TIMESTAMP DEFAULT CURRENT_TIMESTAMP
);
\end{lstlisting}

\subsection{Intégrité des données}

\subsubsection{Contraintes d'intégrité}
\begin{itemize}
  \item \textbf{PRIMARY KEY} : Identifiants uniques
  \item \textbf{UNIQUE} : Email sans doublon
  \item \textbf{FOREIGN KEY} : Références parent-enfant (famille → enfant)
  \item \textbf{CHECK} : Validations (e.g. âge $\geq$ 0)
  \item \textbf{NOT NULL} : Champs obligatoires
\end{itemize}

\subsubsection{Transactions ACID}

Les opérations critiques (inscription avec paiement) utilisent \texttt{BEGIN TRANSACTION} :

\begin{lstlisting}[language=SQL]
BEGIN;
INSERT INTO inscription (famille_id, forfait_id) VALUES (123, 5);
INSERT INTO inscription_detail (inscription_id, ...) VALUES (...);
UPDATE forfait SET places_restantes = places_restantes - 1;
COMMIT;
\end{lstlisting}

Ceci garantit que l'inscription soit atomique (tout ou rien) et évite les doublons lors d'accès concurrents.

% =====================================================
\section{Sécurité}

\subsection{Authentification et Autorisation}

\subsubsection{JWT (JSON Web Token)}
\begin{itemize}
  \item Payload : \texttt{\{id, email, role, iat, exp\}}
  \item Signature : HMAC-SHA256 avec clé secrète serveur
  \item Durée de vie : 24 heures
  \item Revocation : Blacklist en base de données (table \texttt{token\_revoked})
\end{itemize}

\subsubsection{Mot de passe}
\begin{itemize}
  \item Hachage : bcrypt avec salt 10 rounds
  \item Jamais stocké en clair en base de données
  \item Validation : regex email + longueur minimale 12 caractères
\end{itemize}

\subsection{Protection contre les attaques}

\begin{tabular}{lll}
\toprule
\textbf{Menace} & \textbf{Protection} & \textbf{Implémentation} \\
\midrule
SQL Injection & Prepared Statements & \texttt{parameterized queries} \\
XSS & Content Security Policy & Headers HTTP \\
CSRF & CSRF Token + SameSite & Middleware \texttt{csrf.ts} \\
Brute Force & Rate Limiting & Middleware rate-limit \\
\bottomrule
\end{tabular}

\subsection{RGPD et Données personnelles}

\begin{itemize}
  \item Les données des enfants et responsables sont chiffrées en transit (HTTPS/TLS)
  \item Droit à l'oubli : endpoint \texttt{DELETE /api/families/:id} supprime les données
  \item Durée de conservation : 3 ans après dernière inscription
  \item Accord explicite : Checkbox RGPD au moment de l'inscription
\end{itemize}

% =====================================================
\section{Flux Métier Clés}

\subsection{Flux 1 : Inscription en ligne}

\begin{enumerate}
  \item Client accède \texttt{/inscription}
  \item Formulaire réactif (Vue 3) avec validation côté client
  \item \texttt{POST /api/inscriptions} avec données (famille, enfants, forfait)
  \item Serveur valide à nouveau côté serveur (validation défensive)
  \item Transaction atomique : création famille + enfants + inscription
  \item Réponse avec confirmation et numéro de dossier
  \item Email de confirmation envoyé (via service SMTP)
\end{enumerate}

\subsection{Flux 2 : Gestion du planning}

\begin{enumerate}
  \item Admin crée une séance : \texttt{POST /api/sessions}
  \item Données : salle, horaire, ceinture (couleur minimum requise)
  \item Clients consultent le planning via \texttt{GET /api/sessions/public}
  \item Les séances sont filtrées par ceinture (sécurité)
  \item Présence enregistrée : \texttt{PATCH /api/sessions/:id/attendance}
\end{enumerate}

% =====================================================
\section{Déploiement et DevOps}

\subsection{Infrastructure}

\begin{itemize}
  \item \textbf{VPS} : Serveur Linux (Ubuntu 22.04)
  \item \textbf{Node.js} : Runtime Nitro en tant que service systemd
  \item \textbf{PostgreSQL} : DBMS en conteneur Docker ou installé natif
  \item \textbf{Nginx} : Reverse proxy, SSL/TLS, gestion des domaines
  \item \textbf{DNS} : Domaine \texttt{uspv-judo.fr} pointant vers VPS
\end{itemize}

\subsection{Processus CI/CD}

\begin{lstlisting}[language=bash]
# 1. Push sur GitHub
git push origin main

# 2. GitHub Actions trigger
# - Lint (ESLint, Prettier)
# - Unit tests (Vitest)
# - Build (nuxi build)
# - Deploy (SSH vers VPS)

# 3. Sur VPS :
# - Pull du code
# - npm install
# - npm run build
# - Restart du service Node.js
# - Tests de santé (curl /api/health)
\end{lstlisting}

\subsection{Monitoring}

\begin{itemize}
  \item Logs centralisés : stdout → systemd journal
  \item Alertes : crash du service Node.js via systemd
  \item Health check : endpoint \texttt{GET /api/health} vérifie la connexion DB
  \item Backup BD : Script cron \texttt{pg\_dump} quotidien
\end{itemize}

% =====================================================
\section{Décisions Architecturales}

\subsection{Pourquoi Nuxt 3 + Nitro ?}

\begin{tabular}{lp{8cm}}
\toprule
\textbf{Critère} & \textbf{Justification} \\
\midrule
Framework fullstack & Client + API dans un même projet \\
Routing auto & Directories pages/ et api/ auto-compilées \\
SSR optionnel & Rendu serveur pour SEO (futur) \\
TypeScript natif & Type-safety sur client et serveur \\
Écosystème Vue & Composants + stores Pinia cohérents \\
\bottomrule
\end{tabular}

\subsection{Pourquoi PostgreSQL plutôt que NoSQL ?}

\begin{itemize}
  \item \textbf{Relations complexes} : Familles → Enfants → Inscriptions → Forfaits
  \item \textbf{Intégrité référentielle} : Les contraintes FOREIGN KEY evitent les données orphelines
  \item \textbf{Transactions ACID} : Critères pour inscriptions (atomicité)
  \item \textbf{Requêtes complexes} : JOINs sur plusieurs tables (rapports)
\end{itemize}

\subsection{Authentification sans session serveur}

\begin{itemize}
  \item \textbf{JWT stateless} : Pas de stockage de session en mémoire/Redis
  \item \textbf{Scalabilité} : Chaque serveur peut valider le JWT indépendamment
  \item \textbf{Revocation} : Table \texttt{token\_revoked} pour déconnexions urgentes
  \item \textbf{Alternative} : Sessions en Redis (si besoin de contrôle fin ultérieurement)
\end{itemize}

% =====================================================
\section{Points Remarquables et Optimisations}

\subsection{Gestion des erreurs}

\begin{lstlisting}[language=bash]
# Endpoint sécurisé
POST /api/inscriptions
  - Validation côté client (feedback immédiat)
  - Validation côté serveur (défensive)
  - Try/catch autour de la transaction
  - Rollback automatique en cas d'erreur
  - Réponse HTTP 400/500 avec message d'erreur structuré
\end{lstlisting}

\subsection{Performance}

\begin{itemize}
  \item \textbf{Indexes PostgreSQL} : Sur email (unique), id famille (FK), dates (range)
  \item \textbf{Compression} : Gzip des réponses API et assets statiques
  \item \textbf{Cache client} : Immuabilité des assets (cache-busting par hash)
  \item \textbf{Lazy loading} : Composants Nuxt chargés à la demande
\end{itemize}

\subsection{Maintenance et Documentation}

\begin{itemize}
  \item Code source commenté (pourquoi, pas quoi)
  \item Schéma BD documenté (diagramme ERD)
  \item README avec instructions de déploiement
  \item Changelog Git (commits explicites)
\end{itemize}

% =====================================================
\section{Contraintes et Limitations}

\begin{itemize}
  \item \textbf{VPS unique} : Pas de haute disponibilité (pas de cluster)
  \item \textbf{Stockage fichiers} : Pas encore de gestion d'image (avatar, scan diploma)
  \item \textbf{Notifications} : Email uniquement, pas de SMS ni push
  \item \textbf{Statistiques} : Pas de BI ou analytics avancée
\end{itemize}

% =====================================================
\section{Évolutions Futures}

\begin{enumerate}
  \item \textbf{API mobile} : Application iOS/Android avec authentification JWT partagée
  \item \textbf{Paiement en ligne} : Intégration Stripe / PayPal pour règlement forfaits
  \item \textbf{Gestion de fichiers} : Upload avatar, certificat médical, scan identité
  \item \textbf{Notations} : Système d'évaluation des judokas (suivi progression)
  \item \textbf{Clustering} : Réplication BD, load balancer, haute disponibilité
  \item \textbf{Analytics} : Dashboard admin avec rapports d'inscription, taux de présence
\end{enumerate}

% =====================================================
\section{Conclusion}

L'architecture du site USPV Judo repose sur une stack moderne et robuste :
\begin{itemize}
  \item \textbf{Frontend} : SPA réactive (Nuxt 3 / Vue 3)
  \item \textbf{Backend} : API RESTful sécurisée (Nitro / JWT)
  \item \textbf{Données} : Modèle relationnel cohérent (PostgreSQL)
  \item \textbf{Déploiement} : CI/CD automatisé (GitHub Actions)
\end{itemize}

Cette architecture garantit la \textbf{scalabilité}, la \textbf{sécurité} et la \textbf{maintenabilité} du système pour accueillir la croissance future du club USPV Judo.

\vspace{2cm}

\begin{flushright}
  \textit{Mehdi LAFAY}\\
  \textit{Stage BUT2 Informatique}\\
  \textit{IUT Belfort-Montbéliard}\\
  \textit{Juin 2026}
\end{flushright}

\end{document}
